% \iffalse meta-comment
%
% File: lora.dtx
% Copyright 2026 by Dmitriy Nosachev
%
% It may be distributed and/or modified under the conditions of the
% LaTeX Project Public License (LPPL), either version 1.3c of this
% license or (at your option) any later version.  The latest version
% of this license is in the file
%
%    https://www.latex-project.org/lppl.txt
%
%<*driver>
\documentclass{l3doc}
\usepackage{lora}
\usepackage[style=numeric, backend=biber]{biblatex}
\addbibresource{lora.bib}
\usepackage{metalogo}
\usepackage[dvipsnames]{xcolor}
\hypersetup{
  linkcolor=MidnightBlue,
  citecolor=OliveGreen,
  urlcolor=NavyBlue,
  filecolor=Magenta
}
\begin{document}
  \RecordChanges
  \DocInput{\jobname.dtx}
\end{document}
%</driver>
% \fi
% \GetFileInfo{lora.sty}
% \title{\pkg{lora}---\LaTeX{} support of Lora font\\ by Olga Karpushina and Alexei Vanyashin}
% \date{\fileversion, \filedate}
% \author{Dmitriy Nosachev\thanks{\href{mailto:quartz64@gmail.com}{quartz64@gmail.com}}}
% \maketitle
% \begin{abstract}
%   Lora is a well-balanced contemporary serif typeface with roots in calligraphy, featuring brushed curves and driving wedge serifs.
%   This package provides \LaTeX{} support for it.
% \end{abstract}
% \tableofcontents
% \begin{documentation}
%
%\section{User manual}
%\label{sec:ug}
%
% 
%\subsection{Introduction}
%\label{sec:ug-intro}
%
% Lora is a contemporary serif typeface with a balanced design and
% calligraphic origins. This Unicode font family accommodates
% languages using the Latin and Cyrillic scripts (including their
% variants) and has the potential to be extended to support additional
% writing systems.
%
% Lora was created by Olga Karpushina and Alexei Vanyashin at Cyreal
% and launched in 2011 with contributions and assistance from
% Gayaneh Bagdasaryan. Font is licensed under the SIL Open Font
% License, Version 1.1.~\cite{cyreal}
%
% This package provides \LaTeX{} interface for the font.
% Since the font it in OTF format, you do need a Unicode
% engine like \XeTeX{} or \LuaTeX{} to use it.  
%
% I would like to express my deep gratitude to Boris Veytsman,
% the author of the \texttt{arsenal} package~\cite{arsenal}. I used his \texttt{.dtx}
% file (Documented LaTeX Source) as a template.
%
%\subsection{Package options}
%\label{sec:ug-options}
%
% \begin{variable}{default, scale, Scale}
%
% The following options are recognized:
%   \begin{description}
%   \item[default] whether to make Lora the main font of the
%     document, either |true| (the default) or |false|.
%   \item[scale] the scale for the font, by default 1.0.  The option
%   |Scale| is a synonym.
%   \end{description}
% \end{variable}
%
%
%\subsection{Font features}
%\label{sec:ug-families}
%
% \begin{function}{\lorafamily, \textlora}
%   \begin{syntax}
%     \cs{lorafamily} \meta{text}
%     \cs{textlora}\Arg{text}
%   \end{syntax}
%   The font provides the commands \cs{lorafamily} and
%   \cs{textlora} for selecting the font.  Alternatively, the NFSS
%   commands \cs{fontfamily}|{lora}|\cs{selectfont} can be used to
%   select Lora family.
% \end{function}
%
% The font has normal and \textit{italic} shapes, as well as bolded
% \textbf{bold} and \textbf{\itshape bold italic}. They are selected
% by the standard \LaTeX{} commands. 
% 
% This font supports variable weight with values from 400 to 700.
% The regular weight corresponds to 400, and the bold corresponds to 700.
% Additionally, the package provides commands for medium weight (500)
% and semi-bold (600).
% 
% \begin{function}{\LoraMedium, \textloramed}
%   \begin{syntax}
%     \cs{LoraMedium} \meta{text}
%     \cs{textloramed}\Arg{text}
%   \end{syntax}
%   The font provides the commands \cs{LoraMedium} and
%   \cs{textloramed} for selecting the \textloramed{medium weight} variant.
% \end{function}
% 
% \begin{function}{\LoraSemiBold, \textlorasb}
%   \begin{syntax}
%     \cs{LoraSemiBold} \meta{text}
%     \cs{textlorasb}\Arg{text}
%   \end{syntax}
%   The font provides the commands \cs{LoraSemiBold} and
%   \cs{textlorasb} for selecting the \textlorasb{semi-bold} variant.
% \end{function}
%
%\subsection{Special symbols in text}
%\label{sec:chars}
%
% \begin{function}{
%   \texthryvnia,
%   \texttugrik,
%   \texttenge,
%   \textlira,
%   \textruble
% }
% The font has common currency characters, like \cs{textdollar}
% (\textdollar), \cs{textyen} (\textyen), \cs{textsterling}
% (\textsterling), \cs{texteuro} (\texteuro).  It also defines several
% less common currency characters:
%   \cs{texthryvnia} (\texthryvnia),
%   \cs{texttugrik} (\texttugrik),
%   \cs{texttenge} (\texttenge),
%   \cs{textlira} (\textlira),\\
%   \cs{textruble} (\textruble).
% \end{function}
%
% \end{documentation}
% \clearpage
%
% \begin{implementation}
%
% \section{Implementation}
% \label{sec:impl}
%
%\subsection{Setting up}
%\label{sec:settingup}
%
% First, we declare who we are:
%    \begin{macrocode}
%<@@=lora>
%<*package>
\ProvidesExplPackage {lora}
{2026-09-10} {v1.0}
{Lora font by Olga Karpushina and Alexei Vanyashin}
%    \end{macrocode}
%
%\subsection{Options}
%\label{sec:options}
%
% \begin{variable}{
%   default,
%   scale,
%   Scale,
%   \l_@@_default_bool,
%   \l_@@_scale_tl,
% }
%    \begin{macrocode}
\keys_define:nn {lora}

{
  default .bool_set:N = \l__lora_default_bool,
  default .default:n = true,
  scale .tl_set:N = \l__lora_scale_tl,
  Scale .tl_set:N = \l__lora_scale_tl,
}

\keys_set:nn { lora }
{
  default=true,
  scale = 1.0,
}
%    \end{macrocode}
% \end{variable}
%
% Processing options
%    \begin{macrocode}
\IfFormatAtLeastTF { 2022-06-01 }
  { \ProcessKeyOptions [ lora ] }
  {
    \RequirePackage { l3keys2e }
    \ProcessKeysOptions { lora }
  }
%    \end{macrocode}
%
%\subsection{Setting up font}
%\label{sec:font}
%
%    \begin{macrocode}
\RequirePackage{fontspec}

\defaultfontfeatures{Ligatures=TeX, RawFeature={+axis={wght=400}}}
%    \end{macrocode}
% \changes{v1.0.0}{2026/09/10}{Switched to the original font filenames (Lora[wght].ttf, Lora-Italic[wght].ttf).}
%    \begin{macrocode}
\newfontfamily\lorafamily{Lora[wght]}
[
  NFSSFamily         = lora,
  Ligatures          = TeX,
  Scale              = \l__lora_scale_tl,
  Extension          = .ttf,
  ItalicFont         = Lora-Italic[wght],
  BoldFont           = Lora[wght],
  BoldItalicFont     = Lora-Italic[wght],
  BoldFeatures       = {RawFeature={+axis={wght=700}}},
  BoldItalicFeatures = {RawFeature={+axis={wght=700}}}
]

\newfontfamily\LoraMedium{Lora[wght]}[
  Extension          = .ttf,
  UprightFeatures    = {RawFeature={+axis={wght=500}}},
  ItalicFeatures     = {RawFeature={+axis={wght=500}}},
  ItalicFont         = Lora-Italic[wght],
  BoldFont           = Lora[wght],
  BoldItalicFont     = Lora-Italic[wght],
  BoldFeatures       = {RawFeature={+axis={wght=700}}},
  BoldItalicFeatures = {RawFeature={+axis={wght=700}}}
]

\newfontfamily\LoraSemiBold{Lora[wght]}[
  Extension          = .ttf,
  UprightFeatures    = {RawFeature={+axis={wght=600}}},
  ItalicFeatures     = {RawFeature={+axis={wght=600}}},
  ItalicFont         = Lora-Italic[wght],
  BoldFont           = Lora[wght],
  BoldItalicFont     = Lora-Italic[wght],
  BoldFeatures       = {RawFeature={+axis={wght=700}}},
  BoldItalicFeatures = {RawFeature={+axis={wght=700}}}
]
%    \end{macrocode}
%
% Checking whether we want the font to be default
%    \begin{macrocode}
\bool_if:NT \l__lora_default_bool
{
  \setmainfont{lora}
  \renewcommand\rmdefault{lora}
}
%    \end{macrocode}
%
% \begin{macro}{\textlora, \textloramed, \textlorasb}
%    \begin{macrocode}
\DeclareTextFontCommand{\textlora}{\lorafamily}
\DeclareTextFontCommand{\textloramed}{\LoraMedium}
\DeclareTextFontCommand{\textlorasb}{\LoraSemiBold}
%    \end{macrocode}
% \end{macro}
%
% Special characters, absent in the default
% \begin{macro}{
%   \texthryvnia,
%   \texttugrik,
%   \texttenge,
%   \textlira,
%   \textruble
% }
% Currency symbols
%    \begin{macrocode}
\DeclareUnicodeSymbol{\texthryvnia}       {"20B4}
\DeclareUnicodeSymbol{\texttugrik}        {"20AE}
\DeclareUnicodeSymbol{\texttenge}         {"20B8}
\DeclareUnicodeSymbol{\textlira}          {"20BA}
\DeclareUnicodeSymbol{\textruble}         {"20BD}
%    \end{macrocode}
%   
% \end{macro}
%
%    \begin{macrocode}
%</package>
%    \end{macrocode}
% \end{implementation}
%
% \nocite{*}
% \printbibliography
%
%
%
% \PrintIndex
% \PrintChanges
